\documentclass{zhvt-classic}

% 使用本地源流明体（GenRyuMin2 TW）按文件路径加载，不依赖系统字体名
% 版框（内框／界栏／书口侧栏／鱼尾）改用与版心印章同一印泥棕色：
%   印章图片 yshs.png 实色 = RGB(135,68,52) = HTML 874434
\usepackage{xcolor}
\usepackage{tikz}
% sealbrown 双注册：xcolor 供 TikZ/pgf，L3 供 cls 版框
\definecolor{sealbrown}{HTML}{874434}
\ExplSyntaxOn
\color_set:nnn { sealbrown } { HTML } { 874434 }
\ExplSyntaxOff

\zhvtset{
  grid color       = sealbrown,
  main font        = GenRyuMin2TW-SB,
  main font path   = C:/Users/Administrator/texmf/fonts/opentype/public/GenRyuMin2TW/,
jiazhu font      = GenRyuMin2TW-M,
  jiazhu font path = C:/Users/Administrator/texmf/fonts/opentype/public/GenRyuMin2TW/,
  % 页码框容纳 4 字：正常页「九十七」+前置空 1 格=4 行恰好；
  % 100+ 页「一百零一」共 5 行超出，末字「一/二/三」被框裁掉（2026-10-04 实测），故放宽到 5
  shukou page number chars = 5,
  % 书口书名框高从默认 6 字减到 4 字，把省出的 2 字高拨给 chapter 框，
  % 使 6 字书名+卷名放进 9 字高的 chapter 框后不再压线（2026-10-05）
  shukou title chars = 4,
}

\title{史记}[樊郦滕灌列传]
\author{汉·司马迁}
\maker{闲叟漫记·重製}


% ===== 索隐述赞样式（2026-11-08 四次定稿 · 愚大人亲手改定参数）=====
% ⚠ 竖排前提：自定义字体家族必须带 RawFeature={vertical:+vert}；
%   simhei.ttf 加 +vhal 会整体加载失败静默回退，只写 +vert；simfang 同理。
\newCJKfontfamily{\shuzanfont}[Path=C:/Windows/Fonts/,RawFeature={vertical:+vert}]{simhei.ttf}
\newCJKfontfamily{\shuzanbody}[Path=C:/Windows/Fonts/,RawFeature={vertical:+vert}]{simfang.ttf}
% 题头：152x30pt 圆角 6pt 棕底白字；baseline (0,2) 使整框落入界行格子不压版框线；
% 白字节点 x 不补偿、y +8.2pt（像素级墨迹实测 dx/dy≈0，L/R 留白约 20pt）。
\newcommand{\shuzanhead}[1]{%
  \begin{tikzpicture}[baseline={(0,2)}]
    \fill[sealbrown, rounded corners=6pt]
      (-76pt,-15pt) rectangle (76pt,15pt);
    \node[text=white, inner sep=0pt, anchor=center] at (0pt,8.2pt)
      {\shuzanfont\fontsize{23}{23}\selectfont #1};
  \end{tikzpicture}%
}
% ★ 修正书口卷名位置：chapter 框起始位置改回 cls 默认值，避免全局影响。
%   只在每个卷首页通过 \juan 命令里的空命令微调该页位置（安全，不会跨页污染）。
\ExplSyntaxOn
\AtBeginDocument{
  \dim_set:Nn \l__zhvt_shukou_chapter_title_start_dim
    { \l__zhvt_shukou_title_height_dim + \l__zhvt_font_size_dim * 2 }
  \dim_set:Nn \l__zhvt_shukou_page_number_start_dim
    { \l__zhvt_shukou_chapter_title_start_dim + \l__zhvt_shukou_chapter_title_height_dim }
}
\ExplSyntaxOff

\begin{document}

% ★★ 版心（书口）排布覆写：把书名放到鱼尾之下（2026-10-04）
%
%   cls 默认版心自上而下：书名（title 静态框，在鱼尾之上）→ 鱼尾 → 篇名（chapter 动态框）→ 页码
%   本模板改为：          鱼尾（位置不动）→ 书名 → 篇名 → 页码分割线 → 页码
%
%   手法：把 cls 放在 title 静态框里的书名，改放进 chapter 动态框（鱼尾之下）；
%         title 静态框留空。书名取 \l__zhvt_title_str，即 \title{...} 的第一参数，
%         所以换书时只改 \title，不必在此重复写书名。
%
%   ⚠ 三条必须遵守的写法（2026-10-04 实测）：
%     ① 汉字必须用 \hbox{...} 包住。flowfram 会把框内容装进一个宽 14.4mm（= 一个 line_sep
%        = 书口栏宽）的纵向 minipage 并按该宽度折行，而整页靠 PDF /Rotate 90 显示、
%        框内文字本身不变向 —— 不包 \hbox 时「围炉夜话」4 字会被逐字折成 4 行，
%        在版面上就成了横排 4 字并横向溢出版框。
%     ② 竖排下**一个 \hbox 就是一整列**。书名与篇名要同处一列，必须放进同一个 \hbox；
%        分成两个 \hbox 时，后一个会被当成新的一列、甩到隔壁栏去
%        （实测「序」跑到书口栏外侧 14.4mm 的正文列上）。中间的 \quad 在框内是横向空白，
%        显示出来正好是书名与篇名之间的纵向间隔。
%     ③ ★ 内容盒要写成 **与框等宽的定宽盒** `\hbox to \l__zhvt_shukou_chapter_title_height_dim{…\hfil}`。
%        实测：chapter 框是把内容盒沿纵向**居中**摆放的，内容宽度一变，书名就跟着漂：
%        自然宽度的 \hbox（书名 38.4mm）落在 x=157.2mm，加篇名后（57.6mm）落在 x=152.4mm，
%        目录页（无篇名）与正文页的书名就差半格。做成等宽盒后，无论有没有篇名，
%        内容固定从框顶起排，实测两页均落在 x=142.8mm，不再漂。
\ExplSyntaxOn
% 把书口 chapter 框（书名+篇名）起始位置下移 1 字高，避免 4 字卷名顶到鱼尾
% ⚠ cls 的 \__zhvt_calculate_dimensions: 在 \AtBeginDocument 里执行，会重算该 dim；
%   必须把覆写挂到 begin document 之后，否则会被二次覆盖（2026-10-05 实测）。
% (旧方案已移到导言区，直接修改 chapter 框起始位置)
% 卷标题命令 #1=书口篇名(鱼尾) #2=正文卷名（含卷序，卷名内空格用 ~）
%
% ★★ 目录页两项修订（2026-10-05，愚大人指出）★★
%   ① 页码起点统一对齐：
%      原状 —— 目录项写成 \par\hspace{4em}卷名\hspace{4em}页码，页码起点 = 卷名实宽 + 4em。
%      卷名实宽随字数变：「卷一~梁惠王上」6 字 +~、「卷七~离娄上」5 字 +~，
%      差 1 个字位 → 卷七~卷十 4 条的页码整体比其余 10 条高 27.21pt（一个字），看着上下不齐。
%      改法 —— 卷名补全角空格（宽 = 1 字 = 27.215pt）补足到 **7 字符位**再写进目录，
%      14 条页码起点遂全部落在同一条竖线上（起点对齐；末字随页码字数自然参差）。
%      ⚠ 只对目录里的副本 \l_tmpc_tl 补，正文卷名仍用 \l_tmpa_tl 原样 —— 正文像素不受影响。
%   ② 卷名整体上移半字位：前导 \hspace 由 4em 改 3.5em。
%      实测原状目录块在版框内偏下（天头留白 139.6pt / 地脚留白 110.7pt），
%      上移 0.5em = 13.6pt 后上下基本均分，视觉居中。
\NewDocumentCommand { \juan } { mm } {
	\clearpage
	\tl_set:Nn \l__zhvt_chapter_title_tl { #1 }
	\tl_set:Nn \l_tmpa_tl { #2 }
	\tl_set_eq:NN \l_tmpc_tl \l_tmpa_tl
	\int_step_inline:nn { 7 - \str_count:n { #2 } }
	  { \tl_put_right:Nn \l_tmpc_tl { 　 } }
	\tl_set:Nx \l_tmpb_tl { \zhnum { page } }
	\tl_put_left:Nn \l_tmpb_tl { \hspace { 4em } }
	\tl_put_left:Nn \l_tmpb_tl { \l_tmpc_tl }
	\tl_put_left:Nn \l_tmpb_tl { \par \hspace { 3.5em } }
	\addtocontents { toc } { \l_tmpb_tl }
	\skip_horizontal:n { \l__zhvt_font_size_dim * 4 }
	\l_tmpa_tl
	\par
	\thispagestyle { empty }
}
\cs_set_protected:Npn \__zhvt_set_shukou_content: {
	% ★ 鱼尾上方留白（title 静态框）嵌入闲叟钤印：
	%   整页 /Rotate 90，故此处把图用 \rotatebox{90} 转正为竖章；
	%   尺寸卡在版心条幅内（书内约 7mm 宽 × 33.8mm 高，留余量不压鱼尾）。
	\setstaticcontents * { title } {
		\vbox_to_ht:nn { \l__zhvt_baseline_skip_dim } {
			\vskip 14.2pt
			\hbox_to_wd:nn { \l__zhvt_shukou_title_height_dim } {
				\hfil
				\rotatebox [ origin = c ] { 90 } { \includegraphics [ width = 7mm , height = 33.8mm ] { yshs.png } }
				\hfil
			}
		}
	}
% ★ 版心（书口）统一改法（2026-11-16 愚大人定）：18pt＋框内居中＋书口文字同版框色
%   ① 原 27.2pt 时 7 字长篇名（如「廉颇蔺相如列传」）末字溢出 chapter 框、压住页码；
%      改 \fontsize{18pt}{18pt}，前后 \hfil 令字串在框内居中（不随篇名字数漂移）。
%   ② \vbox{\vskip 8.2pt} 定量下移，抵消缩字号后栏宽方向 8.2pt 上偏（像素实测 y 中心 543.3 = 原版）。
%      ⚠ \raisebox / \vspace* 在 flowfram 框内无效；换字号须同步换补偿：16pt→10.1、19pt→7.4。
%   ③ chapter 与 page~number 均置 \color{sealbrown}（HTML 874434），与版框／界行同色；正文仍黑。
%   ⚠ 页码字位、钤印、鱼尾、界行、正文一律未动；五部合订本同此口径。
	\setdynamiccontents * { chapter }
	{ \color{sealbrown} \vbox { \vskip 8.2pt \hbox to \l__zhvt_shukou_chapter_title_height_dim
		{ \hfil \fontsize{18pt}{18pt}\selectfont \l__zhvt_title_str \quad \l__zhvt_chapter_title_tl \hfil } } }
	\setdynamiccontents * { page~number } { \color{sealbrown} \quad \zhnum { page } }
}
\__zhvt_set_shukou_content:
\ExplSyntaxOff

\maketitle

% 单篇，无多卷目录，不渲染目录页

\juan{樊郦滕灌列传}{卷九十五~樊郦滕灌列传第三十五}
\hspace*{2em}舞阳侯樊哙者。沛人也。以屠狗为事。与高祖俱隐。\\
\hspace*{2em}初从高祖起丰。攻下沛。高祖为沛公。以哙为舍人。从攻胡陵。方与。还守丰。击泗水监丰下。破之。复东定沛。破泗水守薛西。与司马枿战砀东。卻敌。斩首十五级。赐爵国大夫。常从。沛公击章邯军濮阳。攻城先登。斩首二十三级。赐爵列大夫。复常从。从攻城阳。先登。下户牖。破李由军。斩首十六级。赐上间爵。从攻围东郡守尉於成武。卻敌。斩首十四级。捕虏十一人。赐爵五大夫。从击秦军。出亳南。河间守军於杠里。破之。击破赵贲军开封北。以卻敌先登。斩候一人。首六十八级。捕虏二十七人。赐爵卿。从攻破杨熊军於曲遇。攻宛陵。先登。斩首八级。捕虏四十四人。赐爵封号贤成君。从攻长社。轘辕。绝河津。东攻秦军於尸。南攻秦军於犨。破南阳守齮於阳城。东攻宛城。先登。西至郦。以卻敌。斩首二十四级。捕虏四十人。赐重封。攻武关。至霸上。斩都尉一人。首十级。捕虏百四十六人。降卒二千九百人。\\
\hspace*{2em}项羽在戏下。欲攻沛公。沛公从百馀骑因项伯面见项羽。谢无有闭关事。项羽既飨军士。中酒。亚父谋欲杀沛公。令项庄拔剑舞坐中。欲击沛公。项伯常蔽之。时独沛公与张良得入坐。樊哙在营外。闻事急。乃持铁盾入到营。营卫止哙。哙直撞入。立帐下。项羽目之。问为谁。张良曰。﹁沛公参乘樊哙。﹂项羽曰。﹁壮士。﹂赐之卮酒彘肩。哙既饮酒。拔剑切肉食。尽之。项羽曰。﹁能复饮乎。﹂哙曰。﹁臣死且不辞。岂特卮酒乎。且沛公先入定咸阳。暴师霸上。以待大王。大王今日至。听小人之言。与沛公有隙。臣恐天下解。心疑大王也。﹂项羽默然。沛公如厕。麾樊哙去。既出。沛公留车骑。独骑一马。与樊哙等四人步从。从间道山下归走霸上军。而使张良谢项羽。项羽亦因遂已。无诛沛公之心矣。是日微樊哙奔入营谯让项羽。沛公事几殆。\\
\hspace*{2em}明日。项羽入屠咸阳。立沛公为汉王。汉王赐哙爵为列侯。号临武侯。迁为郎中。从入汉中。\\
\hspace*{2em}还定三秦。别击西丞白水北。雍轻车骑於雍南。破之。从攻雍。斄城。先登击章平军好畤。攻城。先登陷阵。斩县令丞各一人。首十一级。虏二十人。迁郎中骑将。从击秦车骑壤东。卻敌。迁为将军。攻赵贲。下郿。槐里。柳中。咸阳。灌废丘。最。至栎阳。赐食邑杜之樊乡。从攻项籍。屠煮枣。击破王武。程处军於外黄。攻邹。鲁。瑕丘。薛。项羽败汉王於彭城。尽复取鲁。梁地。哙还至荥阳。益食平阴二千户。以将军守广武。一岁。项羽引而东。从高祖击项籍。下阳夏。虏楚周将军卒四千人。围项籍於陈。大破之。屠胡陵。\\
\hspace*{2em}项籍既死。汉王为帝。以哙坚守战有功。益食八百户。从高帝攻反燕王臧荼。虏荼。定燕地。楚王韩信反。哙从至陈。取信。定楚。更赐爵列侯。与诸侯剖符。世世勿绝。食舞阳。号为舞阳侯。除前所食。以将军从高祖攻反韩王信於代。自霍人以往至云中。与绛侯等共定之。益食千五百户。因击陈豨与曼丘臣军。战襄国。破柏人。先登。降定清河。常山凡二十七县。残东垣。迁为左丞相。破得綦毋卬。尹潘军於无终。广昌。破豨别将胡人王黄军於代南。因击韩信军於参合。军所将卒斩韩信。破豨胡骑横谷。斩将军赵既。虏代丞相冯梁。守孙奋。大将王黄。将军。太仆解福等十人。与诸将共定代乡邑七十三。其後燕王卢绾反。哙以相国击卢绾。破其丞相抵蓟南。定燕地。凡县十八。乡邑五十一。益食邑千三百户。定食舞阳五千四百户。从。斩首百七十六级。虏二百八十八人。别。破军七。下城五。定郡六。县五十二。得丞相一人。将军十二人。二千石已下至三百石十一人。\\
\hspace*{2em}哙以吕后女弟吕须为妇。生子伉。故其比诸将最亲。\\
\hspace*{2em}先黥布反时。高祖尝病甚。恶见人。卧禁中。诏户者无得入群臣。群臣绛。灌等莫敢入。十馀日。哙乃排闼直入。大臣随之。上独枕一宦者卧。哙等见上流涕曰。﹁始陛下与臣等起丰沛。定天下。何其壮也。今天下已定。又何惫也。且陛下病甚。大臣震恐。不见臣等计事。顾独与一宦者绝乎。且陛下独不见赵高之事乎。﹂高帝笑而起。\\
\hspace*{2em}其後卢绾反。高帝使哙以相国击燕。是时高帝病甚。人有恶哙党於吕氏。即上一日宫车晏驾。则哙欲以兵尽诛灭戚氏。赵王如意之属。高帝闻之大怒。乃使陈平载绛侯代将。而即军中斩哙。陈平畏吕后。执哙诣长安。至则高祖已崩。吕后释哙。使复爵邑。\\
\hspace*{2em}孝惠六年。樊哙卒。谥为武侯。子伉代侯。而伉母吕须亦为临光侯。高后时用事专权。大臣尽畏之。伉代侯九岁。高后崩。大臣诛诸吕。吕须婘属。因诛伉。舞阳侯中绝数月。孝文帝既立。乃复封哙他庶子市人为舞阳侯。复故爵邑。市人立二十九岁卒。谥为荒侯。子他广代侯。六岁。侯家舍人得罪他广。怨之。乃上书曰。﹁荒侯市人病不能为人。令其夫人与其弟乱而生他广。他广实非荒侯子。不当代後。﹂诏下吏。孝景中六年。他广夺侯为庶人。国除。\\
\hspace*{2em}曲周侯郦商者。高阳人。陈胜起时。商聚少年东西略人。得数千。沛公略地至陈留。六月馀。商以将卒四千人属沛公於岐。从攻长社。先登。赐爵封信成君。从沛公攻缑氏。绝河津。破秦军洛阳东。从攻下宛。穰。定十七县。别将攻旬关。定汉中。\\
\hspace*{2em}项羽灭秦。立沛公为汉王。汉王赐商爵信成君。以将军为陇西都尉。别将定北地。上郡。破雍将军焉氏。周类军栒邑。苏駔军於泥阳。赐食邑武成六千户。以陇西都尉从击项籍军五月。出钜野。与锺离眛战。疾斗。受梁相国印。益食邑四千户。以梁相国将从击项羽二岁三月。攻胡陵。\\
\hspace*{2em}项羽既已死。汉王为帝。其秋。燕王臧荼反。商以将军从击荼。战龙脱。先登陷阵。破荼军易下。卻敌。迁为右丞相。赐爵列侯。与诸侯剖符。世世勿绝。食邑涿五千户。号曰涿侯。以右丞相别定上谷。因攻代。受赵相国印。以右丞相赵相国别与绛侯等定代。雁门。得代丞相程纵。守相郭同。将军已下至六百石十九人。还。以将军为太上皇卫一岁七月。以右丞相击陈豨。残东垣。又以右丞相从高帝击黥布。攻其前拒。陷两陈。得以破布军。更食曲周五千一百户。除前所食。凡别破军三。降定郡六。县七十三。得丞相。守相。大将各一人。小将二人。二千石已下至六百石十九人。\\
\hspace*{2em}商事孝惠。高后时。商病。不治。其子寄。字况。与吕禄善。及高后崩。大臣欲诛诸吕。吕禄为将军。军於北军。太尉勃不得入北军。於是乃使人劫郦商。令其子况绐吕禄。吕禄信之。故与出游。而太尉勃乃得入据北军。遂诛诸吕。是岁商卒。谥为景侯。子寄代侯。天下称郦况卖交也。\\
\hspace*{2em}孝景前三年。吴。楚。齐。赵反。上以寄为将军。围赵城。十月不能下。得俞侯栾布自平齐来。乃下赵城。灭赵。王自杀。除国。孝景中二年。寄欲取平原君为夫人。景帝怒。下寄吏。有罪。夺侯。景帝乃以商他子坚封为缪侯。续郦氏後。缪靖侯卒。子康侯遂成立。遂成卒。子怀侯世宗立。世宗卒。子侯终根立。为太常。坐法。国除。\\
\hspace*{2em}汝阴侯夏侯婴。沛人也。为沛厩司御。每送使客还。过沛泗上亭。与高祖语。未尝不移日也。婴已而试补县吏。与高祖相爱。高祖戏而伤婴。人有告高祖。高祖时为亭长。重坐伤人。告故不伤婴。婴证之。後狱覆。婴坐高祖系岁馀。掠笞数百。终以是脱高祖。\\
\hspace*{2em}高祖之初与徒属欲攻沛也。婴时以县令史为高祖使。上降沛一日。高祖为沛公。赐婴爵七大夫。以为太仆。从攻胡陵。婴与萧何降泗水监平。平以胡陵降。赐婴爵五大夫。从击秦军砀东。攻济阳。下户牖。破李由军雍丘下。以兵车趣攻战疾。赐爵执帛。常以太仆奉车从击章邯军东阿。濮阳下。以兵车趣攻战疾。破之。赐爵执珪。复常奉车从击赵贲军开封。杨熊军曲遇。婴从捕虏六十八人。降卒八百五十人。得印一匮。因复常奉车从击秦军雒阳东。以兵车趣攻战疾。赐爵封转为滕公。因复奉车从攻南阳。战於蓝田。芷阳。以兵车趣攻战疾。至霸上。项羽至。灭秦。立沛公为汉王。汉王赐婴爵列侯。号昭平侯。复为太仆。从入蜀。汉。\\
\hspace*{2em}还定三秦。从击项籍。至彭城。项羽大破汉军。汉王败。不利。驰去。见孝惠。鲁元。载之。汉王急。马罢。虏在後。常蹶两兒欲弃之。婴常收。竟载之。徐行面雍树乃驰。汉王怒。行欲斩婴者十馀。卒得脱。而致孝惠。鲁元於丰。\\
\hspace*{2em}汉王既至荥阳。收散兵。复振。赐婴食祈阳。复常奉车从击项籍。追至陈。卒定楚。至鲁。益食兹氏。\\
\hspace*{2em}汉王立为帝。其秋。燕王臧荼反。婴以太仆从击荼。明年。从至陈。取楚王信。更食汝阴。剖符世世勿绝。以太仆从击代。至武泉。云中。益食千户。因从击韩信军胡骑晋阳旁。大破之。追北至平城。为胡所围。七日不得通。高帝使使厚遗阏氏。冒顿开围一角。高帝出欲驰。婴固徐行。弩皆持满外向。卒得脱。益食婴细阳千户。复以太仆从击胡骑句注北。大破之。以太仆击胡骑平城南。三陷陈。功为多。赐所夺邑五百户。以太仆击陈豨。黥布军。陷陈卻敌。益食千户。定食汝阴六千九百户。除前所食。\\
\hspace*{2em}婴自上初起沛。常为太仆。竟高祖崩。以太仆事孝惠。孝惠帝及高后德婴之脱孝惠。鲁元於下邑之间也。乃赐婴县北第第一。曰﹁近我﹂。以尊异之。孝惠帝崩。以太仆事高后。高后崩。代王之来。婴以太仆与东牟侯入清宫。废少帝。以天子法驾迎代王代邸。与大臣共立为孝文皇帝。复为太仆。八岁卒。谥为文侯。子夷侯灶立。七年卒。子共侯赐立。三十一年卒。子侯颇尚平阳公主。立十九岁。元鼎二年。坐与父御婢奸罪。自杀。国除。\\
\hspace*{2em}颍阴侯灌婴者。睢阳贩缯者也。高祖之为沛公。略地至雍丘下。章邯败杀项梁。而沛公还军於砀。婴初以中涓从击破东郡尉於成武及秦军於扛里。疾斗。赐爵七大夫。从攻秦军亳南。开封。曲遇。战疾力。赐爵执帛。号宣陵君。从攻阳武以西至雒阳。破秦军尸北。北绝河津。南破南阳守齮阳城东。遂定南阳郡。西入武关。战於蓝田。疾力。至霸上。赐爵执珪。号昌文君。\\
\hspace*{2em}沛公立为汉王。拜婴为郎中。从入汉中。十月。拜为中谒者。从还定三秦。下栎阳。降塞王。还围章邯於废丘。未拔。从东出临晋关。击降殷王。定其地。击项羽将龙且。魏相项他军定陶南。疾战。破之。赐婴爵列侯。号昌文侯。食杜平乡。\\
\hspace*{2em}复以中谒者从降下砀。以至彭城。项羽击。大破汉王。汉王遁而西。婴从还。军於雍丘。王武。魏公申徒反。从击破之。攻下黄。西收兵。军於荥阳。楚骑来众。汉王乃择军中可为骑将者。皆推故秦骑士重泉人李必。骆甲习骑兵。今为校尉。可为骑将。汉王欲拜之。必。甲曰。﹁臣故秦民。恐军不信臣。臣原得大王左右善骑者傅之。﹂灌婴虽少。然数力战。乃拜灌婴为中大夫。令李必。骆甲为左右校尉。将郎中骑兵击楚骑於荥阳东。大破之。受诏别击楚军後。绝其饷道。起阳武至襄邑。击项羽之将项冠於鲁下。破之。所将卒斩右司马。骑将各一人。击破柘公王武。军於燕西。所将卒斩楼烦将五人。连尹一人。击王武别将桓婴白马下。破之。所将卒斩都尉一人。以骑渡河南。送汉王到雒阳。使北迎相国韩信军於邯郸。还至敖仓。婴迁为御史大夫。\\
\hspace*{2em}三年。以列侯食邑杜平乡。以御史大夫受诏将郎中骑兵东属相国韩信。击破齐军於历下。所将卒虏车骑将军华毋伤及将吏四十六人。降下临菑。得齐守相田光。追齐相田横至嬴。博。破其骑。所将卒斩骑将一人。生得骑将四人。攻下嬴。博。破齐将军田吸於千乘。所将卒斩吸。东从韩信攻龙且。留公旋於高密。卒斩龙且。生得右司马。连尹各一人。楼烦将十人。身生得亚将周兰。\\
\hspace*{2em}齐地已定。韩信自立为齐王。使婴别将击楚将公杲於鲁北。破之。转南。破薛郡长。身虏骑将一人。攻阳。前至下相以东南僮。取虑。徐。度淮。尽降其城邑。至广陵。项羽使项声。薛公。郯公复定淮北。婴度淮北。击破项声。郯公下邳。斩薛公。下下邳。击破楚骑於平阳。遂降彭城。虏柱国项佗。降留。薛。沛。酂。萧。相。攻苦。谯。复得亚将周兰。与汉王会颐乡。从击项籍军於陈下。破之。所将卒斩楼烦将二人。虏骑将八人。赐益食邑二千五百户。\\
\hspace*{2em}项籍败垓下去也。婴以御史大夫受诏将车骑别追项籍至东城。破之。所将卒五人共斩项籍。皆赐爵列侯。降左右司马各一人。卒万二千人。尽得其军将吏。下东城。历阳。渡江。破吴郡长吴下。得吴守。遂定吴。豫章。会稽郡。还定淮北。凡五十二县。\\
\hspace*{2em}汉王立为皇帝。赐益婴邑三千户。其秋。以车骑将军从击破燕王臧荼。明年。从至陈。取楚王信。还。剖符。世世勿绝。食颍阴二千五百户。号曰颍阴侯。\\
\hspace*{2em}以车骑将军从击反韩王信於代。至马邑。受诏别降楼烦以北六县。斩代左相。破胡骑於武泉北。复从击韩信胡骑晋阳下。所将卒斩胡白题将一人。受诏并将燕。赵。齐。梁。楚车骑。击破胡骑於硰石。至平城。为胡所围。从还军东垣。\\
\hspace*{2em}从击陈豨。受诏别攻豨丞相侯敞军曲逆下。破之。卒斩敞及特将五人。降曲逆。卢奴。上曲阳。安国。安平。攻下东垣。\\
\hspace*{2em}黥布反。以车骑将军先出。攻布别将於相。破之。斩亚将楼烦将三人。又进击破布上柱国军及大司马军。又进破布别将肥诛。婴身生得左司马一人。所将卒斩其小将十人。追北至淮上。益食二千五百户。布已破。高帝归。定令婴食颖阴五千户。除前所食邑。凡从得二千石二人。别破军十六。降城四十六。定国一。郡二。县五十二。得将军二人。柱国。相国各一人。二千石十人。\\
\hspace*{2em}婴自破布归。高帝崩。婴以列侯事孝惠帝及吕太后。太后崩。吕禄等以赵王自置为将军。军长安。为乱。齐哀王闻之。举兵西。且入诛不当为王者。上将军吕禄等闻之。乃遣婴为大将。将军往击之。婴行至荥阳。乃与绛侯等谋。因屯兵荥阳。风齐王以诛吕氏事。齐兵止不前。绛侯等既诛诸吕。齐王罢兵归。婴亦罢兵自荥阳归。与绛侯。陈平共立代王为孝文皇帝。孝文皇帝於是益封婴三千户。赐黄金千斤。拜为太尉。\\
\hspace*{2em}三岁。绛侯勃免相就国。婴为丞相。罢太尉官。是岁。匈奴大入北地。上郡。令丞相婴将骑八万五千往击匈奴。匈奴去。济北王反。诏乃罢婴之兵。後岁馀。婴以丞相卒。谥曰懿侯。子平侯阿代侯。二十八年卒。子彊代侯。十三年。彊有罪。绝二岁。元光三年。天子封灌婴孙贤为临汝侯。续灌氏後。八岁。坐行赇有罪。国除。\\
\hspace*{2em}太史公曰。吾適丰沛。问其遗老。观故萧。曹。樊哙。滕公之家。及其素。异哉所闻。方其鼓刀屠狗卖缯之时。岂自知附骥之尾。垂名汉廷。德流子孙哉。余与他广通。为言高祖功臣之兴时若此云。\\

% ===== 索隐述赞（司马贞《史记索隐》）· 统一样式 =====
% 题头「索隐述赞曰」五字并入反白框；正文每四句一列，定宽盒 \hfil 居列中；
% \fontsize{24}{41} 列距对齐全书界行（121+41k）；\clearpage 使述赞独占一页。
\clearpage
\shuzanhead{索隐述赞曰}\linebreak
{\shuzanbody\fontsize{24}{41}\selectfont
\hbox to \hsize{\hfil 圣贤影响。云蒸龙变。屠狗贩缯。攻城野战。\hfil}\linebreak
\hbox to \hsize{\hfil 扶义西上。受封南面。郦况卖交。舞阳内援。\hfil}\linebreak
\hbox to \hsize{\hfil 滕灌更王。奕叶繁衍。\hfil}%
\par}
\end{document}
